\section{The remaining quadratic-eight label at coefficient three}
\label{sec:coefficient-three-negative8-strict}

\paragraph{A different local argument.}
The exclusion of $\chi_{-8}$ in G178 was a limitation of its local
invariant argument.  Its value at three is one, so that argument cannot
be applied to this label.  We instead prove that one decomposition
group generates the maximal pro-three quotient over its quadratic
splitting field.  The proof uses explicit global units and a single
$S$-unit, on the original coefficient ring throughout.

Put $K=\mathbf Q(\omega)$, $r=\sqrt{-2}$, $t=\sqrt{-3}=1+2\omega$,
and $L=K(r)$.  For a finite extension $E/\mathbf Q_3$ with its original
ring of integers $O$ and maximal ideal $\mathfrak m$, let
\[
 D=(E/O)(\chi_{-8}\kappa^{-1}c_3^{-e}),
 \qquad e\geq1,\qquad
 \kappa:G_K\longrightarrow1+\mathfrak m
 \quad\hbox{unramified away three}.
 \tag{G187.1}
\]
The strict conditions mean zero in full $H^1$ at the unique place of
$K$ over three and zero in $H^1(I_x,D)$ at every other finite place.

\begin{theorem}[AC-G187: the negative quadratic-eight label]
\label{AC-G187}
For every original $O$, every positive $e$ and every $\kappa$ in
(G187.1), the full strict group is zero:
\[
                         S_3(K,D)=0.                 \tag{G187.2}
\]
The same holds on the actual group $G_{K,\Sigma_3}$ of (FC.2), with
every original local and inertia condition.  In particular the pure
finite label $\chi_{-8}$ is supplied at coefficient three for every
positive Tate weight, without any condition on the ramification of $O$.
No conductor-seven, nontrivial cubic or auxiliary factor is added.
\end{theorem}

\begin{proof}
\textbf{1. The small class-group input.}
The two quadratic fields $K$ and $\mathbf Q(r)$ are distinct; hence
$[L:\mathbf Q]=4$ and $L$ is totally imaginary with two complex pairs.
The product order $\mathbf Z[\omega,r]$ has discriminant
\[
                  (-3)^2(-8)^2=576.                  \tag{G187.3}
\]
Indeed the two trace matrices have determinants $-3$ and $-8$;
the product trace matrix is their tensor product.  The field
discriminant therefore has absolute value at most $576$, without
asserting that this order is maximal.  The ideal-class representative
bound (Milne, \emph{Algebraic Number Theory}, version 3.08, Theorem 4.3)
is at most
\[
 \left(\frac4\pi\right)^2\frac{4!}{4^4}\sqrt{576}
                    =\frac{36}{\pi^2}<4.             \tag{G187.4}
\]
Consequently every ideal class has an integral representative of norm
at most three.  At two the residue degree is at least two because
$X^2+X+1$ is irreducible modulo two in the subfield $K$; hence no prime
of $L$ has norm two.

The polynomial $X^2+2$ has two simple roots modulo three, lifting to
$\mathbf Q_3$.  Write $v,w$ for the two places of $L$ above three,
chosen so that $r\equiv-1\pmod3$ at $v$ and $r\equiv1\pmod3$ at $w$.
Both completions are $\mathbf Q_3(t)$, of ramification degree two and
residue degree one.  Put $\alpha=1+r$.  Since
$(1+r)(1-r)=3$, the element $\alpha$ has valuation zero at $w$ and
valuation two at $v$.  Its rational norm in $L$ is $9$, so there are
no other prime factors:
\[
                  (\alpha)=v^2,\qquad(1-r)=w^2.
                                                        \tag{G187.5}
\]
The only possible nontrivial norm-at-most-three representatives are
$v$ and $w$, each of class order dividing two.  Thus every ideal class
has order dividing two and, in particular,
\[
                             \operatorname{Cl}(L)[3]=0.
                                                        \tag{G187.6}
\]
No class-group computation or fundamental-unit assertion is used.

\textbf{2. Three explicit generators of the local unit cube quotient.}
At $w$ put $F=\mathbf Q_3(t)$ and $R=\mathbf Z_3[t]$.
The polynomial $X^2+3$ is Eisenstein, so $R$ is its integer ring and
$t$ is a uniformizer.  Set $U_n=1+t^nR$.  Its residue field is
$\mathbf F_3$, and $\omega=(-1+t)/2$ has order three with
$\omega-1$ of valuation one.  Therefore
\[
 R^\times=\{\pm1\}U_1,\qquad U_1=\mu_3U_2,
       \qquad U_1^3=U_2^3=U_4.                        \tag{G187.7}
\]
For the last equality use $U_2=1+3R$, $U_4=1+9R$ and
\[
        (1+3c)^3=1+9(c+3c^2+3c^3).
\]
Every $b\in R$ has a unique solution of
$c+3c^2+3c^3=b$: iteration of $c\mapsto b-3c^2-3c^3$ is a
contraction on the complete ring $R$.  This proves the equality
directly, with no logarithmic-lattice premise.
Multiplication identifies
\[
                  U_2/U_4\simeq 3R/9R
                     =\mathbf F_3[3]\oplus\mathbf F_3[3t].
                                                        \tag{G187.8}
\]
Here $1+x$ corresponds to $x$, and products of two elements of $3R$
vanish modulo $9R$.

The chosen root $r$ at $w$ satisfies $r\equiv4\pmod9$, since
$r\equiv1\pmod3$ and $r^2=-2$.  Thus
$-\alpha=-(1+r)\equiv4\pmod9$, and its class in (G187.8) is $3$.
The global unit
\[
                     \epsilon=5+2rt=5+2\sqrt6
\]
has inverse $5-2rt$.  Direct multiplication gives
\[
 \begin{split}
 u=(-\epsilon)\omega
   &=\frac{5+6r+(-5+2r)t}{2},\\
 u-1&=\frac{3(1+2r)}2+\frac{(2r-5)t}{2}
           \equiv\frac32t\pmod{9R}.
 \end{split}                                            \tag{G187.9}
\]
The constant term is in $27\mathbf Z_3$, and $2r-5\equiv3\pmod9$.
Consequently $u\in U_2$, and its class supplies a nonzero multiple
of the $3t$ direction.  The classes of $-\alpha$ and $u$ span
$U_2/U_4$, while $\omega$ supplies $U_1/U_2$.  Since $-1$ is itself
a cube, (G187.7)--(G187.9) prove that
\[
 R^\times/(R^\times)^3
       \quad\hbox{is generated by }\omega,\epsilon,\alpha.
                                                        \tag{G187.10}
\]
Although $\alpha$ is only a global $S$-unit, it is a unit at this
particular place $w$; its support at the other place is exactly (G187.5).

\textbf{3. No cyclic cubic quotient splits at the chosen place.}
Let $M/L$ be a cyclic extension of degree three, unramified away
$\{v,w\}$, in which $v$ splits completely.  For each place $x$ let
$\rho_x:L_x^\times\to\operatorname{Gal}(M/L)$ be its local Artin map,
with image in the decomposition group.  The product of these maps
is one on $L^\times$ by global reciprocity.  At every place other
than $v,w$, units map to one because the extension is unramified;
the complex places have no nontrivial finite extension.  At $v$ the
entire map is one because $v$ splits.
Applying the product formula to the global units $\omega,\epsilon$
therefore shows $\rho_w(\omega)=\rho_w(\epsilon)=1$.
The same formula applies to $\alpha$: it is a unit at every place
outside $v$, and the whole map at $v$ is trivial.  Thus
$\rho_w(\alpha)=1$ as well.  The target has exponent three, so
(G187.10) forces $\rho_w(R^\times)=1$.
The local reciprocity identification of units with inertia shows
that $M/L$ is also unramified at $w$.  It is now everywhere unramified,
contrary to (G187.6) and Hilbert class-field reciprocity.  Hence such
an $M$ does not exist.

Let $\Gamma=G_{L,\{v,w\}}$ and let $P=\Gamma^{(3)}$ be its maximal
pro-three quotient.  The image $H$ of $G_{L_v}$ in $P$ is compact,
hence closed.  If $H\ne P$, an open normal subgroup of $P$ separates
some element of $P$ from $H$.  In the resulting finite three-group,
the proper image of $H$ lies in a maximal subgroup, which is normal
of index three.  Its quotient would give precisely the excluded
extension $M/L$.  Therefore
\[
                         G_{L_v}\longrightarrow P
                         \quad\hbox{is surjective}.    \tag{G187.11}
\]
The finite-group argument is the one proved in G182, Steps 1--2;
only its arithmetic input has changed.  Joint generation by the
two local groups would not suffice, and is not what is claimed here.

\textbf{4. The local coboundary is a global one on the original ring.}
Over $L$, the character $\chi_{-8}$ is trivial.  For example,
$r=\zeta_8+\zeta_8^3$ has square $-2$, and the actions of
$1,3,5,7$ modulo eight on it are $1,1,-1,-1$, exactly $\chi_{-8}$.
Since $L$ contains $\mu_3$, the remaining action
$a=\kappa^{-1}c_3^{-e}$ takes values in $1+\mathfrak m$ and is
unramified outside $\{v,w\}$.

Let $z:G_L\to D$ be the restriction of a strict cocycle over $K$.
At every inertia group away three the coefficient over $L$ is
trivial, so the local coboundary condition says that the actual
cocycle is zero there.  Consequently $z$ descends to $\Gamma$.
At $v$ there is $d\in D$ such that $z(g)=(a(g)-1)d$.
Choose $m$ so large that the finite module $D[\varpi^m]$ contains
both $z(\Gamma)$ and $d$, where $\varpi$ is a uniformizer of the
original $O$.  Compactness gives finiteness of the cocycle image.
The homomorphism
\[
 \Gamma\longrightarrow
 D[\varpi^m]\rtimes\operatorname{im}(a|_{D[\varpi^m]}),
            \qquad g\longmapsto(z(g),a(g))             \tag{G187.12}
\]
has finite three-group image: both the additive module and the
near-one scalar image are three-groups.  It factors through $P$.
By (G187.11), every global element has the same image as some
local element of $G_{L_v}$.  The local coboundary identity therefore
holds for every $g\in\Gamma$, with the same $d$.  This proves that
the restriction of the strict class to $G_L$ is zero.

Restriction from $K$ to $L$ is injective on $H^1(K,D)$ because
$[L:K]=2$ is invertible on the three-primary module $D$; equivalently,
restriction followed by corestriction is multiplication by two.
Thus the original class is zero.  This argument retains the original
$E/O$ and does not use vanishing or divisibility of a local invariant
group.  It proves (G187.2).

\textbf{5. Actual growing-host coefficients.}
For $h\geq1$, the original host (FG.3) contains $\mathbf Q(\zeta_8)$,
so $\chi_{-8}$ is one actual finite label.  On its original completion
$O_3$ the scalar identification is exactly
\[
 (O_3\otimes_{\mathbf Z_3}\mathbf Q_3/\mathbf Z_3(\chi_{-8}))(-e)
       =(\operatorname{Frac}(O_3)/O_3)(\chi_{-8}c_3^{-e}).
                                                        \tag{G187.13}
\]
Inflation from the actual group in (FC.2) preserves the full local
condition at three and every inside-set inertia condition.  Outside
the defining ramification set the inflated cocycle is already zero
on inertia.  If the inflated cocycle is a global coboundary, its
witness gives the same coboundary on the quotient.  Thus the actual
strict group injects into the zero group just proved.  All extra
ramified primes of the host are retained, and the conclusion supplies
one new actual label at each of the $T-1$ weights $1\leq e<T$.
\end{proof}

\paragraph{Sources and boundary.}
The class bound is the precise Milne Theorem 4.3 already read with
its full applied proof for G185; no whole-book reading is claimed.
The retained local and global class-field theorems are cited at
Sharifi, \emph{Algebraic Number Theory}, Definition 10.4.4,
Theorem 10.4.9(a), and Theorem 9.4.7(b), whose full printed proof
identifies local units with inertia.  The exact sources are
\url{https://www.math.ucla.edu/~sharifi/notes/algnum-ch10.html}
and \url{https://www.math.ucla.edu/~sharifi/notes/algnum-ch09.html};
Hilbert class-field reciprocity is Sharifi, \emph{Iwasawa Theory},
Theorem 1.1.11.  Their precise applied statements are read and
recorded; the deep foundations of class field theory are retained.
The finite group and semidirect cocycle arguments were proved in
G182; the new explicit local cube quotient is proved above.
The geometric interpretation is the retained finite quadratic cover
over $\operatorname{Spec}O_K[1/6]$ and its original rank-one
Artin--Tate sheaf, with the previously audited SGA 1, Expos\'e V,
sections 4 and 7, dictionary.  No new Grothendieck archive attribution
or mixed motive is claimed.  A separate same-window comparison is
still needed before changing the numerical remainder; neither an
individual positive derived-index zero, the full all-prime estimate,
$\mathcal R_{\mathrm{cub}}$, nor H04.4 follows from this strict theorem
alone.  No Lean build or new class computation is claimed.
